% chapters/bab4-hasil-dan-pembahasan.tex -- BAB 4 HASIL DAN PEMBAHASAN
% Contoh teks ilmiah di bawah ini adalah placeholder untuk memperlihatkan
% hasil format (thesis.md); ganti dengan isi penelitian yang sebenarnya.

\chapter{HASIL DAN PEMBAHASAN}

\section{Hasil Praproses Data}
Tahapan praproses data menghasilkan korpus ulasan produk yang telah
dibersihkan dari karakter tidak relevan, dinormalisasi dari kata
tidak baku, dan dihilangkan kata hentinya (stopword). Setelah tahapan
praproses, jumlah token unik pada korpus berkurang secara signifikan
dibandingkan dengan korpus mentah, yang menunjukkan bahwa normalisasi
kata tidak baku berhasil menyatukan variasi ejaan informal ke dalam
bentuk kata baku yang konsisten. Contoh hasil praproses untuk satu
ulasan ditampilkan pada Gambar~\ref{fig:praproses}.

\begin{figure}[htbp]
  \centering
  \fbox{\parbox[c][3.5cm][c]{0.8\linewidth}{\centering
    Contoh diagram tahapan praproses teks\\
    (ganti dengan \texttt{\textbackslash includegraphics})}}
  \caption{Contoh hasil tahapan praproses data ulasan}
  \label{fig:praproses}
\end{figure}

\section{Hasil Pelatihan Model}
Model LSTM dilatih menggunakan data latih selama sejumlah epoch
hingga nilai fungsi kerugian (\textit{loss}) pada data validasi tidak
lagi menurun secara signifikan. Grafik nilai \textit{loss} dan
akurasi selama proses pelatihan menunjukkan bahwa model mulai
menunjukkan gejala \textit{overfitting} setelah epoch tertentu,
sehingga digunakan mekanisme \textit{early stopping} untuk menghentikan
pelatihan pada titik kinerja validasi terbaik.

\section{Hasil Pengujian}
Kinerja model LSTM yang diusulkan dibandingkan dengan model
\textit{baseline} berbasis SVM pada data uji. Ringkasan hasil
pengujian disajikan pada Tabel~\ref{tab:hasil-pengujian}.

\begin{table}[htbp]
  \centering
  \caption{Perbandingan kinerja model pada data uji}
  \label{tab:hasil-pengujian}
  \singlespaced{%
    \begin{tabular}{lcccc}
      \hline
      Model    & Akurasi & Presisi & Recall & F1-score \\
      \hline
      SVM (baseline) & 0.78 & 0.76 & 0.75 & 0.75 \\
      LSTM (usulan)  & 0.87 & 0.86 & 0.85 & 0.85 \\
      \hline
    \end{tabular}%
  }
\end{table}

Berdasarkan Tabel~\ref{tab:hasil-pengujian}, model LSTM yang diusulkan
memperoleh nilai F1-score yang lebih tinggi dibandingkan model
\textit{baseline} pada seluruh kelas sentimen, dengan peningkatan
yang paling terlihat pada kelas netral.

\section{Pembahasan}
Peningkatan kinerja model LSTM dibandingkan dengan model SVM diduga
disebabkan oleh kemampuan LSTM dalam menangkap hubungan kontekstual
antar kata dalam suatu kalimat, termasuk pola negasi yang sering
menjadi sumber kesalahan klasifikasi pada model berbasis fitur kata
tunggal seperti SVM. Sebagai contoh, kalimat yang mengandung frasa
negasi seperti ``tidak bagus'' cenderung diklasifikasikan secara
tepat oleh model LSTM karena model tersebut mempertimbangkan urutan
kata, sedangkan model SVM yang menggunakan representasi
\textit{bag-of-words} cenderung salah mengklasifikasikan kalimat
tersebut sebagai sentimen positif akibat munculnya kata ``bagus''.

Meskipun demikian, model LSTM masih menunjukkan tingkat kesalahan
yang relatif tinggi pada ulasan yang mengandung sarkasme atau opini
implisit, yang secara umum juga menjadi tantangan pada penelitian
analisis sentimen lainnya \citep{example2022}. Temuan ini menunjukkan
bahwa masih terdapat ruang untuk perbaikan pada penelitian
selanjutnya, misalnya dengan memanfaatkan model berbasis Transformer
yang memiliki kemampuan pemahaman konteks yang lebih baik.
